% Validated Results staging fragment for ICNP 2026 venue draft.
% This file contains only owner-validated reductions and uses only packages
% already loaded by ICNP_2026_venue_draft.tex.

\section{Results}
\label{sec:SimulationResults}

Using the matched grid defined in \cref{sec:studydesign}, we report validated master-dataset results by \cref{sec:research_questions}: RQ1 tests stochastic viability, RQ2 tests threat escalation, and RQ3 tests allocator/capacity deployment effects. GA-report observations are used only as supporting context.

% ============================================================================
% RQ1: Stochastic Routing Viability
% ============================================================================
\subsection{RQ1: Stochastic Routing Viability}
\label{subsec:rq1_answer}

\noindent\textbf{Takeaway.}
Under the Stochastic scenario in \cref{tab:scenario_specs}, contextual pursuit and epsilon-greedy baselines remain viable, while several context-free or poorly matched variants collapse. This establishes baseline viability before the structured and adaptive threats analyzed in RQ2.

\smallTitle{RQ1 evidence slice}
RQ1 uses the EXP-family, CMAB, and iCMAB master datasets, filtered to Stochastic runs under the default allocator. Results aggregate across 4K--8K horizons, $s\in\{1,1.5,2\}$, both $T$/$T_b$ replay semantics, and the 3- and 5-run suites.

\begin{table}[ht!]
\footnotesize
\centering
\setlength{\tabcolsep}{4pt}
\caption{RQ1 stochastic results from the validated master datasets. Values are mean Oracle-normalized efficiency.}
\label{tab:rq1_master_stochastic}
\begin{tabular}{p{0.42\linewidth} c c c}
\toprule
\textbf{Model} & \textbf{3 runs} & \textbf{5 runs} & \textbf{Avg.} \\
\midrule
\multicolumn{4}{l}{\textit{Top tier: viable under stochastic disruption}} \\
CPursuit & 89.3 & 90.6 & \textbf{89.9} \\
CEpsilonGreedy & 87.3 & 88.4 & 87.8 \\
iCEpsilonGreedy & 87.3 & 88.4 & 87.8 \\
GNeuralUCB & 83.7 & 87.4 & 85.5 \\
\midrule
\multicolumn{4}{l}{\textit{Mid tier: degraded}} \\
EXPNeuralUCB & 81.6 & 79.6 & 80.6 \\
EXPUCB & 75.1 & 73.9 & 74.5 \\
CThompsonSampling & 65.5 & 67.9 & 66.7 \\
iCPursuit & 66.7 & 68.0 & 67.4 \\
iCThompsonSampling & 60.4 & 61.7 & 61.1 \\
\midrule
\multicolumn{4}{l}{\textit{Collapsed: structural failure}} \\
CEXP4 & 38.2 & 41.0 & 39.6 \\
CEpochGreedy & 35.6 & 38.3 & 37.0 \\
iCEpochGreedy & 35.6 & 38.3 & 37.0 \\
iCEXP4 & 35.6 & 38.3 & 37.0 \\
\bottomrule
\end{tabular}
\end{table}

\Cref{tab:rq1_master_stochastic} shows a sharp stochastic-regime split: CPursuit and epsilon-greedy contextual variants remain at or above the 85\% viability target, while CEXP4 and epoch-greedy variants collapse to roughly 37--40\%. Thus, RQ1 is answered affirmatively: stochastic decoherence alone separates viable contextual baselines from structural failures. RQ3 later disaggregates the capacity effects hidden inside this aggregate.

% ============================================================================
% RQ2: Robustness Under Adaptive Threats
% ============================================================================
\subsection{RQ2: Robustness Under Adaptive Threats}
\label{subsec:rq2_answer}

\noindent\textbf{Takeaway.}
Under Markov, Adaptive, and OnlineAdaptive threats, contextual and informed policies maintain stronger robustness floors than EXP-family adversarial-first baselines. This extends RQ1 from stochastic viability to adaptive disruption.

\smallTitle{RQ2 evidence slice}
RQ2 uses the Markov, Adaptive, and OnlineAdaptive threat slice from \cref{tab:scenario_specs}, restricted to the default allocator and aggregated across 4K--8K horizons, $s\in\{1,1.5,2\}$, both $T$/$T_b$ replay semantics, and the 3- and 5-run suites. The table reports the strongest CMAB and iCMAB representatives together with two EXP-family adversarial baselines.

\begin{table}[ht!]
\footnotesize
\centering
\setlength{\tabcolsep}{4pt}
\caption{RQ2 adversarial-scope results. Win Dominance is each displayed model's share of aggregated scenario wins among the four representatives under the locked Markov/Adaptive/OnlineAdaptive scope.}
\label{tab:rq2_adversarial}
\begin{tabular}{p{0.30\linewidth} c c c c}
\toprule
\textbf{Algorithm} & \textbf{Avg.} & \textbf{CV} & \textbf{Floor} & \textbf{Win Dom.} \\
\midrule
CPursuit & \textbf{88.1} & 5.3 & 77.4 & 25.6 \\
iCEpsilonGreedy & 86.9 & \textbf{3.6} & \textbf{81.0} & \textbf{43.9} \\
EXPNeuralUCB & 82.4 & 16.5 & 18.0 & 30.5 \\
EXP3/UCB & 76.3 & 6.0 & 68.8 & 0.0 \\
\bottomrule
\end{tabular}
\end{table}

\Cref{tab:rq2_adversarial} refutes the adversarial-first hypothesis: CPursuit leads average efficiency, iCEpsilonGreedy provides the strongest stability floor, and EXPNeuralUCB remains fragile despite adversarial exploration. Thus, RQ2 is answered affirmatively: structured and adaptive threats expose stability gaps not visible in RQ1, and adversarial-first EXP-family baselines do not dominate the locked adversarial scope.

% Second-look candidates before final venue cleanup: robustness frontier figure (fig:floor) and explicit RQ2a/RQ2b/RQ2c supporting-answer bullets.

% ============================================================================
% RQ3: Deployment Optimization
% ============================================================================
\subsection{RQ3: Deployment Optimization}
\label{subsec:rq3}

\noindent\textbf{Takeaway.}
RQ3 evaluates whether algorithm, allocator, and replay-capacity choices should change with the threat regime. Deployment is configuration-sensitive, but the hybrid dataset also reveals a strong static configuration that exceeds the 85\% robustness target across all five scenarios.

\smallTitle{RQ3 evidence slice}
RQ3 uses the validated Hybrid master dataset for CPursuitNeuralUCB, iCPursuitNeuralUCB, GNeuralUCB, and EXPNeuralUCB across all scenarios, allocators, replay anchoring types, scales $s\in\{1,1.5,2\}$, and 4K--8K horizons. Unless otherwise noted, point estimates average the 3- and 5-run suites; RQ3 CV values are computed across scenario means as a deployment-volatility proxy because per-run variance is not available in the master table.

\smallTitle{Validated RQ3 answer}
Deployment is configuration-sensitive, but the Hybrid master dataset also reveals a strong static configuration. At the 6K horizon, \texttt{iCPursuitNeuralUCB + Fixed + ($T$-type, $s=2$)} achieves a 96.0\% global average and a 92.8\% worst-case floor across scenarios, using the validated mean over 3- and 5-run suites. Scenario-specific optima still differ, so later RQ3 analyses disaggregate allocator and replay-capacity effects.

% Keep for second-look before final venue cleanup: threat-adaptive allocator figure (fig:threat_rules) and capacity--efficiency figure (fig:capacity_all). The latter likely belongs with RQ3b.

% ============================================================================
% RQ3a: Predictive Context Modeling Impact
% ============================================================================
\subsubsection{RQ3a: Predictive Context Modeling Impact}
\label{subsubsec:rq3a_answer}

\noindent\textbf{Focus.}
RQ3a isolates whether the informative/predictive variant, iCPursuitNeuralUCB, improves deployment robustness over reactive CPursuitNeuralUCB.

\smallTitle{RQ3a evidence slice}
To isolate architecture impact, we hold deployment fixed at 6K horizon, \texttt{Fixed} allocator, $T$-type anchoring, and $s=2$, then use the validated 3-run suite from the Hybrid master dataset.

\begin{table}[ht!]
\footnotesize
\centering
\setlength{\tabcolsep}{3pt}
\caption{RQ3a informative-context impact under the fixed deployment setting. Values are from the validated 3-run Hybrid master-dataset provenance branch.}
\label{tab:rq3a_master_informative}
\begin{tabular}{p{0.30\linewidth} c c c c c c c}
\toprule
\textbf{Model} & \textbf{Bl} & \textbf{Sh} & \textbf{Mk} & \textbf{Ag} & \textbf{OA} & \textbf{Avg.} & \textbf{CV} \\
\midrule
CPursuitNeuralUCB & 99.8 & 94.7 & 93.0 & 92.8 & 81.5 & 92.4 & 6.5 \\
iCPursuitNeuralUCB & 99.9 & 94.8 & 93.0 & 92.8 & 99.8 & 96.0 & 3.3 \\
\bottomrule
\end{tabular}

{\footnotesize Bl=Baseline, Sh=Stochastic, Mk=Markov, Ag=Adaptive, OA=OnlineAdaptive; CV is computed across scenario means.}
\end{table}

\Cref{tab:rq3a_master_informative} shows that iCPursuitNeuralUCB improves global robustness mainly by lifting OnlineAdaptive efficiency from 81.5\% to 99.8\% (+18.3 pp) while keeping Stochastic, Markov, and Adaptive performance effectively unchanged. It also reduces scenario-level dispersion from 6.5 to 3.3 CV. Thus, informative context is most valuable when the threat adapts online.

% ============================================================================
% RQ3b: Replay Capacity Scaling and Paradox
% ============================================================================
\subsection{RQ3b: Replay Capacity Scaling \& Paradox}
\label{subsec:rq3b_answer}

\noindent\textbf{Focus.}
RQ3b tests whether larger replay scale improves robustness or creates a threat-dependent capacity paradox.

\smallTitle{RQ3b evidence slice}
RQ3b fixes the allocator to \texttt{Random}, the horizon to 6K, and the replay anchoring to $T$-type, then varies $s\in\{1,1.5,2\}$ using Hybrid master-dataset means over the 3- and 5-run suites. This isolates replay-scale behavior before allocator co-design is analyzed in RQ3c.

\begin{table}[ht!]
\footnotesize
\centering
\setlength{\tabcolsep}{3pt}
\caption{RQ3b capacity scaling under fixed \texttt{Random} allocation, 6K horizon, and $T$-type replay anchoring. Values are Hybrid master-dataset means over the 3- and 5-run suites.}
\label{tab:rq3b_master_capacity_scaling}
\begin{tabular}{p{0.30\linewidth} c c c c c c c}
\toprule
\textbf{Model} & \textbf{$s$} & \textbf{Bl} & \textbf{Sh} & \textbf{Mk} & \textbf{Ag} & \textbf{OA} & \textbf{Avg.} \\
\midrule
CPursuitNeuralUCB & 1.0 & 94.7 & 93.8 & 89.3 & 88.5 & 81.5 & 89.6 \\
                  & 1.5 & 45.8 & 77.1 & 51.5 & 88.5 & 69.0 & 66.4 \\
                  & 2.0 & 96.3 & 94.9 & 95.9 & 88.5 & 84.5 & 92.0 \\
\midrule
iCPursuitNeuralUCB & 1.0 & 92.1 & 93.6 & 88.7 & 88.5 & 70.8 & 86.7 \\
                   & 1.5 & 56.1 & 73.2 & 42.2 & 88.5 & 83.5 & 68.7 \\
                   & 2.0 & 95.1 & 94.5 & 96.3 & 88.5 & 83.5 & 91.6 \\
\bottomrule
\end{tabular}

{\footnotesize Bl=Baseline, Sh=Stochastic, Mk=Markov, Ag=Adaptive, OA=OnlineAdaptive.}
\end{table}

\Cref{tab:rq3b_master_capacity_scaling} shows that replay scale is not a monotonic ``more is better'' knob. Under fixed \texttt{Random} allocation and $T$-type anchoring, both pursuit-based hybrids degrade sharply at $s=1.5$ and recover at $s=2$, confirming that replay capacity must be tuned jointly with allocator and anchoring type. Extended GA ablations remain supporting evidence, but the master-dataset trend is the validated RQ3b conclusion.

% Keep for second-look before final venue cleanup: capacity--efficiency figure (fig:capacity_all), likely with RQ3b if a visual capacity-paradox artifact is needed.
